\section{The Byzantine Line}\label{sec:byzantine}

The Greek Church after Dionysius received his nine orders as the Church's
own speech about heaven and kept the angels at the centre of her theology
and her worship. John Damascene gathered the teaching of the Greek Fathers
into the chapters on the angels and the demons that the Latin schools would
read as the \emph{De fide orthodoxa}; Maximus the Confessor read the Divine
Liturgy as the entry of the faithful into the angels' choir; the councils
and the monks of the iconoclast age confessed that the angels are rightly
painted and honoured; Theodore the Studite preached the angels on their
feast; Andrew of Caesarea read their war and their priestly service in the
Apocalypse; and Gregory Palamas set man's image and the angels' light side
by side. The liturgy of Constantinople, sung in the Byzantine rite to this
day, places the angels in the entrances, the hymns, and the anaphora of
every Divine Liturgy, and keeps their feast on 8 November and the memory of
Michael's miracle at Chonae on 6 September. These witnesses stand beside the
Western line as a second voice of the one tradition; they agree with it on
the angels' nature, ministry, and fall, and they speak in their own idiom on
the time of the angels' creation, the rank of the fallen, the image of God,
and the light in which the angels see.

\sectionguard
\subsection{John Damascene}\label{sec:bz-damascene}

The \emph{Exposition of the Orthodox Faith} is the third part of
Damascene's \emph{Fount of Knowledge}. It is a summary, often in the
Fathers' own words, and its chapters on the angels (II.3) and on the devil
and demons (II.4) gather what the Greek Fathers taught before him. The
Latin West read it in the twelfth-century translation of Burgundio of
Pisa, and Aquinas cites it throughout the treatise on the angels. English quotations are from the translation of S.
D. F. Salmond in the \emph{Nicene and Post-Nicene Fathers}; the Greek is
Migne's (PG~94, 865--880).

\paragraph{The definition of the angel.} Damascene begins with the
Creator. God ``brought them out of nothing into being and created them after
His own image, an incorporeal race, a sort of spirit or immaterial fire,''
as David sings, ``He maketh His angels spirits, and His ministers a flame of
fire'' (Ps 103[104]:4); the image of fire describes ``their lightness and
the ardour, and heat, and keenness and sharpness with which they hunger for
God and serve Him, and how they are borne to the regions above and are
quite delivered from all material thought'' (\emph{De fide orth.}\ II.3).
Then comes the definition:

\begin{quote}
An angel, then, is an intelligent essence, in perpetual motion, with
free-will, incorporeal, ministering to God, having obtained by grace an
immortal nature: and the Creator alone knows the form and limitation of its
essence. But all that we can understand is, that it is incorporeal and
immaterial. For all that is compared with God Who alone is incomparable, we
find to be dense and material. For in reality only the Deity is immaterial
and incorporeal. (II.3)
\end{quote}

\noindent The Greek runs: \gk{Angelos toinyn estin ousia noera,
aeikin\=etos, autexousios, as\=omatos, The\=o leitourgousa, kata charin en
t\=e physei to athanaton eil\=ephuia} (PG~94, 865). Each term carries a
doctrine. \gk{Ousia noera}, an intellectual essence, places the angel in
the order of mind. \gk{Aeikin\=etos}, ever-moving, names the unceasing act
of an intelligence always turned toward God; Aquinas, meeting the word in an
objection that would make the angel a body, explains that ``an angel is
called an ever mobile substance, because he is ever actually intelligent,
and not as if he were sometimes actually and sometimes potentially, as we
are'' (\ST{I q.~50 a.~1 ad~2}). \gk{Autexousios}, self-determining, is
free will. \gk{As\=omatos}, bodiless, is said relatively: compared with God
the angel is ``dense and material,'' a sentence Aquinas reads as the logic
of a mean between extremes, so that angels, compared to God, seem
corporeal, ``not, however, as if anything corporeal existed in them'' (ad~1).
\gk{The\=o leitourgousa}, ministering to God, gives the angel his office in
the definition itself. And the immortality \gk{kata charin}, by grace,
places even the angel's endurance under the gift of God: ``For all that has
had beginning comes also to its natural end. But God alone is eternal.''
Aquinas answers the objection drawn from this phrase by saying that
Damascene ``is dealing with perfect immortality, which includes complete
immutability,'' which the angels have from God (\ST{I q.~50 a.~5 ad~1}).

\paragraph{Freedom, the absence of repentance, and immortality.} ``The
angel's nature then is rational, and intelligent, and endowed with
free-will, changeable in will, or fickle. For all that is created is
changeable, and only that which is un-created is unchangeable. Also all that
is rational is endowed with free-will'' (II.3). Created and free, the angel
could ``abide or progress in goodness, or \dots\ turn towards evil.'' But
once turned it does not turn back: ``It is not susceptible of repentance
because it is incorporeal. For it is owing to the weakness of his body that
man comes to have repentance.'' The good angels, ``With difficulty \dots\
moved to evil, yet \dots\ not absolutely immoveable,'' are now ``altogether
immoveable, not by nature but by grace and by their nearness to the Only
Good.'' The Latin schools teach the same confirmation of the good angels in
beatitude (\ST{I q.~62 a.~8}; see \ref{sec:q62}).

\paragraph{Lights, place, and rank.} The angels ``are secondary intelligent
lights derived from that first light which is without beginning, for they
have the power of illumination; they have no need of tongue or hearing, but
without uttering words they communicate to each other their own thoughts and
counsels.'' They were created through the Word and perfected by the Spirit:
``Through the Word, therefore, all the angels were created, and through the
sanctification by the Holy Spirit were they brought to perfection, sharing
each in proportion to his worth and rank in brightness and grace.'' They are
circumscribed, not by walls or bars but by their createdness: ``when they
are in the Heaven they are not on the earth: and when they are sent by God
down to the earth they do not remain in the Heaven.'' Being minds, ``they
are in mental places,'' present and at work ``after the manner of a mind''
where they are assigned, and they ``cannot be present and energise in
various places at the same time.'' Aquinas takes the first of these
sentences for the \emph{sed contra} of his article on whether an angel can
be in several places at once (\ST{I q.~52 a.~2}; see \ref{sec:q52}).

On whether the angels differ in essence Damascene professes ignorance:
``Whether they are equals in essence or differ from one another we know
not. God, their Creator, Who knoweth all things, alone knoweth. But they
differ from each other in brightness and position, whether it is that their
position is dependent on their brightness, or their brightness on their
position: and they impart brightness to one another, because they excel one
another in rank and nature. And clearly the higher share their brightness
and knowledge with the lower.'' Aquinas, following the logic of pure forms,
determines that each angel is its own species (\ST{I q.~50 a.~4}); the Greek
doctor leaves the question to God and states only the difference of light
and station and the descent of illumination from higher to lower.

\paragraph{The ministry and the life of the angels.} ``They are mighty and
prompt to fulfil the will of the Deity, and their nature is endowed with
such celerity that wherever the Divine glance bids them there they are
straightway found. They are the guardians of the divisions of the earth:
they are set over nations and regions, allotted to them by their Creator:
they govern all our affairs and bring us succour.'' ``They behold God
according to their capacity, and this is their food.'' ``They are above us
for they are incorporeal, and are free of all bodily passion, yet are not
passionless: for the Deity alone is passionless.'' ``They take different
forms at the bidding of their Master, God, and thus reveal themselves to men
and unveil the divine mysteries to them. They have Heaven for their
dwelling-place, and have one duty, to sing God's praise and carry out His
divine will.'' Contemplation, praise, and ministry are one life.

\paragraph{The nine orders.} Damascene gives the nine orders on the
authority of ``that most holy, and sacred, and gifted theologian, Dionysius
the Areopagite,'' who ``divides into three groups, each containing three''
the nine names of Scripture. The Greek names the first triad \gk{t\=en t\=on
hexapteryg\=on Seraphim, kai t\=on polyommat\=on Cheroubim, kai t\=on
hagi\=otat\=on Thron\=on}, the six-winged Seraphim, the many-eyed Cherubim,
and the most holy Thrones, who are ever about God and united to him
immediately; the second, \gk{Kyriot\=et\=on, kai Dyname\=on, kai
Exousi\=on}, the Dominations, Virtues, and Powers; and the third, \gk{Arch\=on,
kai Archangel\=on, kai Angel\=on}, the Principalities, Archangels, and
Angels (PG~94, 873--874). The order is Dionysius's own (\emph{CH} 6.2; see
\ref{sec:ch}); the Greek line does not take up Gregory the Great's
exchange of the Virtues and Principalities (see \ref{sec:hierarchies}).

\paragraph{The time of their creation.} ``Some, indeed, like Gregory the
Theologian, say that these were before the creation of other things. He
thinks that the angelic and heavenly powers were first and that thought was
their function. Others, again, hold that they were created after the first
heaven was made. But all are agreed that it was before the foundation of
man. For myself, I am in harmony with the theologian. For it was fitting
that the mental essence should be the first created, and then that which can
be perceived, and finally man himself, in whose being both parts are
united.'' Aquinas knows this to be the common mind of the Greeks. Answering
Jerome's word about the ages in which the orders served God before the
world, he says: ``Jerome is speaking according to the teaching of the Greek
Fathers; all of whom hold the creation of the angels to have taken place
previously to that of the corporeal world'' (\ST{I q.~61 a.~3 ad~1}). He
judges creation together with the corporeal world ``the more probable,''
because the angels are part of one universe with the bodies, and adds that
``the contrary is not to be deemed erroneous; especially on account of the
opinion of Gregory Nazianzen'' (\emph{ibid.}; see \ref{sec:q61}). Damascene
closes the chapter by excluding the angels from creation's cause: ``those who
say that the angels are creators of any kind of essence whatever are the
mouth of their father, the devil. For since they are created things they
are not creators.''

\paragraph{The devil and the demons (II.4).} The fourth chapter opens with
the fallen prince:

\begin{quote}
He who from among these angelic powers was set over the earthly realm, and
into whose hands God committed the guardianship of the earth, was not made
wicked in nature but was good, and made for good ends, and received from his
Creator no trace whatever of evil in himself. But he did not sustain the
brightness and the honour which the Creator had bestowed on him, and of his
free choice was changed from what was in harmony to what was at variance
with his nature, and became roused against God Who created him, and
determined to rise in rebellion against Him: and he was the first to depart
from good and become evil. (II.4)
\end{quote}

\noindent ``For evil is nothing else than absence of goodness, just as
darkness also is absence of light.'' With him ``an innumerable host of
angels subject to him were torn away and followed him and shared in his
fall.'' Their power is bounded: ``they have no power or strength against any
one except what God in His dispensation hath conceded to them, as for
instance, against Job and those swine that are mentioned in the Gospels.''
They do not know the future, though they guess at it and know the
Scriptures; and ``while the liberty to attack man has been granted to them,
they have not the strength to over-master any one: for we have it in our
power to receive or not to receive the attack.'' The chapter ends with the
sentence the Latin schools repeated: ``what in the case of man is death is a
fall in the case of angels. For after the fall there is no possibility of
repentance for them, just as after death there is for men no repentance''
(\gk{hoti hoper esti tois anthr\=opois ho thanatos, touto tois angelois h\=e
ekpt\=osis}; PG~94, 877--878).

Aquinas uses this chapter at many points. The saying on death and the fall
grounds his account of the obstinacy of the demons (\ST{I q.~64 a.~2}; see
\ref{sec:q64}). The prince ``set over the earthly realm'' is his witness
that the holy doctors, with the Platonists, placed angelic powers over
corporeal things (\ST{I q.~110 a.~1}). He cites Damascene for the view that
all the demons sinned at one time (\ST{I q.~63 a.~8} obj.~1), and he weighs
Damascene's account of the devil's rank against Gregory's. Damascene, he
says, holds that ``the highest of those who sinned was set over the
terrestrial order,'' so that ``they who fell were of the lower grade of
angels''; Gregory holds that ``he who sinned was the very highest of all.''
Aquinas judges Gregory's the more probable view, since pride's motive is
excellence, ``Yet this must not be prejudicial to the other view'' (\ST{I
q.~63 a.~7}; see \ref{sec:q63}). The Greek line keeps Damascene's reading:
Andrew of Caesarea, below, says the devil was entrusted first with the rule
of the air.

\sectionguard
\subsection{Maximus the Confessor: The Liturgy as the Angels' Choir}\label{sec:bz-maximus}

Maximus the Confessor wrote the \emph{Mystagogia} as an explanation of
the Church and of her Liturgy, in which Dionysius is his guide. The angels enter it
at the hymn that the faithful sing with them. The Greek is Migne's
(PG~91), with a rendering after each passage.

In the thirteenth chapter Maximus reads the whole sequence of the rite
after the Gospel as the Word's leading of the faithful to the Father.
Having gathered them by the kiss of peace and the confession of the Creed,
the Word, \gk{eita tois angelois enarithmious autous katast\=esas dia tou
Trisagiou, kai t\=en aut\=en ekeinois epist\=em\=en t\=es hagiastik\=es
theologias autois charisamenos, t\=o The\=o kai Patri prosagei}
(\emph{Myst.}\ 13; PG~91, 692C--D): then, having numbered them with the
angels through the thrice-holy hymn and given them the same knowledge of the
hallowing theology as the angels have, he brings them to God the Father,
adopted in the Spirit through the prayer in which they are made worthy to
call God Father. The hymn of the Seraphim (Isa 6:3) is the step by which men
are counted among the angels before they are counted among the sons.

The nineteenth chapter asks what the doxology of the thrice-holy signifies:

\begin{quote}
\gk{H\=e de ginomen\=e triss\=e tou hagiasmou t\=es theias hymnologias
ekbo\=esis para pantos tou pistou laou, t\=en pros tas as\=omatous kai
noeras dynameis kata to mellon phan\=esomen\=en hen\=osin te kai isotimian
parad\=eloi; kath' h\=en symph\=on\=os tais an\=o dynamesi, dia taut\=ot\=eta
t\=es atreptou peri Theon aeikin\=esias, trisin hagiasmois hymnein te kai
hagiazein t\=en trisypostaton mian theot\=eta didachth\=esetai t\=on
anthr\=op\=on h\=e physis.} (\emph{Myst.}\ 19; PG~91, 696C)
\end{quote}

\noindent The threefold cry of sanctification in the divine hymn, raised by
the whole faithful people, shows the union and equality of honour with the
bodiless and intellectual powers that will be manifested in the age to
come; in it human nature, in concord with the powers above, through the
sameness of an unchanging ever-moving motion about God, will be taught to
hymn and hallow with three sanctifications the one Godhead in three
hypostases. The angels' \gk{aeikin\=esia}, Damascene's ``perpetual motion,''
is here the destiny of the saints. The Liturgy sings now what the
resurrection will make complete.

In the summary of his last chapter Maximus distinguishes what the angelic
hymn means for all and what it means for each. In general it signifies the
equality and common life of the heavenly and earthly powers in the age to
come, when the risen body no longer weighs down the soul; for the believers
it means \gk{t\=en pros angelous kata t\=en pistin theologik\=en hamillan},
a theological emulation of the angels according to faith; for those given
to ascetic practice, a life that rivals the angels' splendour as far as man
may (\emph{Myst.}\ 24; PG~91, 709--710). The monk's ``angelic life'' and the
layman's ``Holy, holy, holy'' are one vocation at two degrees.

\sectionguard
\subsection{The Angels in Image}\label{sec:bz-images}

The iconoclast controversy put a question about the angels to the whole
Church: may they be painted, since they are bodiless? Damascene's three
apologies against those who attack the holy images answer that they may.
``Holy Scripture clothes in figure God and the angels,'' and Dionysius
explains why (\emph{De imag.}\ I; tr.\ Allies, 1898). In the third apology
he sets out the principle:

\begin{quote}
Now if nothing physical or material may be attributed to an angel, a
spirit, and a devil, yet they may be depicted and circumscribed after their
own nature. Being intellectual beings, they are believed to be present and
to energise in places known to us intellectually. They are represented
materially as Moses made an image of the cherubim who were looked upon by
those worthy of the honour, the material image offering them an immaterial
and intellectual sight. Only the divine nature is uncircumscribed and
incapable of being represented in form or shape, and incomprehensible.
(\emph{De imag.}\ III)
\end{quote}

\noindent He repeats the teaching of the \emph{Exposition}: ``The angel,
and a soul, and a demon, compared to God, who alone is incomparable, are
bodies; but compared to material bodies, they are bodiless.'' God, ``not
wishing that we should be in ignorance of spirits, clothed them in type and
figure.'' The honour given to the angels is not the worship of God: ``Josue
saw the image of an angel, not as he is, for an angel is not visible to
bodily eyes, and falling down he adored, and so did Daniel. Yet an angel is
a creature, and servant, and minister of God, not God. And he worshipped the
angel not as God, but as God's ministering spirit'' (\emph{ibid.}; Jos
5:15; Dan 8:17).

The Second Council of Nicaea defined the matter in 787.
\latin{Sequentesque divinitus inspiratum sanctorum Patrum nostrorum
magisterium, et catholicae traditionem Ecclesiae}, following the divinely
inspired teaching of the holy Fathers and the tradition of the Catholic
Church, it decreed that venerable and holy images be set up in
churches, on sacred vessels and vestments, on walls and panels, in houses
and by the roads, \latin{tam videlicet imaginem Domini Dei et Salvatoris
nostri Jesu Christi, quam intemeratae Dominae nostrae sanctae Dei
genitricis, honorabiliumque Angelorum, et omnium Sanctorum} (DS~600): the
image of our Lord and God and Saviour Jesus Christ, of our immaculate Lady
the holy Mother of God, of the honourable angels, and of all the saints.
The council distinguished the honour paid to images, with kisses, incense,
and lights, from \latin{veram latriam, quae secundum fidem est quaeque solam
divinam naturam decet}, the true worship of adoration that belongs to the
divine nature alone, \latin{Imaginis enim honor ad primitivum transit}, for
the honour of the image passes to its original (DS~601). The fourth council
of Constantinople (870), which Denzinger numbers the eighth ecumenical
council, confirmed the practice in its third canon:
\latin{insuper et iconas sanctorum Angelorum depingimus, quemadmodum eos
figurat verbis divina Scriptura} (DS~656)---we also paint the icons of the
holy angels, as divine Scripture portrays them in words.

\sectionguard
\subsection{Theodore the Studite: A Sermon for the Synaxis}\label{sec:bz-theodore}

Among the orations of Theodore the Studite Angelo Mai printed a sermon \gk{Eis t\=en synaxin t\=on ourani\=on tagmat\=on}, for
the synaxis of the heavenly orders, under Theodore's name, noting that an altered excerpt of it had been
printed among the spurious works of Chrysostom (\emph{Or.}\ 6,
\emph{In sanctos angelos}; PG~99, 729--736). In his encomium on the
Dormition Theodore sets the nine orders about the Mother of God (see
\ref{sec:cm-queen}); this sermon is given wholly to the angels.

\paragraph{The festival of the two choirs.} It opens with trumpets: the
trumpets of the spiritual commanders have sounded around us, the graces of
the divine ministers have shone about us, the ranks of wakeful hymn-singers
have gathered us to hymns, and the immaterial orders have joined the
material in one dance. \gk{Ho t\=on an\=o chorostat\=es kai t\=on kat\=o,
en t\=o aut\=o oik\=o taxiarchei s\=emeron} (\emph{Or.}\ 6.1; PG~99,
729): the leader of the choir above and below marshals them today in the
same house. He is Michael, \gk{pros charitos Theou lamprynomenos, kai
brotois the\=on ellamps\=e\=on mesit\=es genomenos}, radiant with the grace
of God and made the mediator of divine illuminations to mortals. The angels
are the patrons and intercessors of men, \gk{philanthr\=opoi angeloi}, who
have received their love of man from the Master who loves mankind; they are
\gk{theoi t\=e the\=osei ek Theou, all' ou t\=e physei legomenoi}, called
gods by deification from God and not by nature; and they are \gk{to
empsychon kai noeron pyr, h\=e logik\=e Theou phlox}, the living and
intellectual fire, the rational flame of God (6.1; PG~99, 732; Ps
103[104]:4).

\paragraph{One essence, many offices.} The second section gives Theodore's
account of the orders: \gk{diaphoros h\=e stasis, diaphoroi hai
pros\=egoriai, diaphoroi hai axiai, mia h\=e ousia, polytropous echousa tas
ergasias} (6.2; PG~99, 732)---different stations, different names, different
dignities, one essence with manifold works. There are orders that sing,
orders that guard the heavens, orders that stand immediately before God;
inner and outer orders, the outer called the gates of heaven; higher and
lower orders, the lower once crying to the higher, ``Lift up your gates, O
ye princes'' (Ps 23[24]:7); orders that send and orders that are sent,
orders that initiate and orders that are initiated. Angels fly through the
air and rescue the souls of the saints from the spirits of the air, as
Jacob said at his death, ``The angel that delivereth me from all evils''
(Gen 48:16); angels are entrusted with the administration of the earthly
world; an angel is given by God to every God-fearing soul, as the Lord
said, ``their angels in heaven always see the face of my Father'' (Matt
18:10); and one angel is assigned to every faithful nation, for God
``set the bounds of the nations according to the number of the angels of
God'' (Deut 32:8, Septuagint, tr.\ Brenton). Then comes a litany of their works:
through angels the mysteries of God are imparted on earth; through angels
waters and springs are guarded and healed; through angels the worthy are
delivered from temptations; through angels the relics of the martyrs are
kept; through angels the altars of God are guarded and the churches of the
faithful protected; through angels petitions and prayers are carried from
earth to heaven and divine gifts given to men. \gk{Kai pan mikrou dein epi
g\=es Theou di' angel\=on teleitai myst\=erion} (6.2; PG~99, 733): almost
every mystery of God on earth is accomplished through angels.

\paragraph{The angels of the Old Law.} The third section walks through
sacred history: the Cherubim set to guard the tree of life; the three angels
at Mamre, \gk{ton trisagion t\=es Triados m\=enyonta thesmon}, disclosing the
thrice-holy ordinance of the Trinity; the angel who wrestled with Jacob; the
flame of angels in the bush that burned and was not consumed, prefiguring
the unconfused union of God with the thorny nature of mortals; the angel who
slew the firstborn of Egypt; the angel in the pillar of fire, guiding the
people as Christ; the many angels at Sinai with trumpets, fire, and cloud;
the angel who announced victory to Josue (6.3; PG~99, 733--736). Theodore
preaches the angels as Damascene defines them: one intelligent nature,
graded in light, sent in every age for the salvation of men.

\sectionguard
\subsection{Andrew of Caesarea on the Apocalypse}\label{sec:bz-andrew}

Andrew of Caesarea wrote a commentary on the Apocalypse in sermons and
chapters, printed in Migne's hundred and sixth volume of the Greek
Fathers. Two of its chapters give the angels their priestly service and
their war.

\paragraph{The angel of the censer.} His twenty-first chapter is headed as the
opening of the seventh seal, showing the angelic powers bringing to God the
prayers of the saints as incense (\emph{In Apoc.}\ 21; PG~106). The text is
Apoc 8:3--4: ``another angel came and stood before the altar, having a
golden censer: and there was given to him much incense, that he should offer
of the prayers of all saints, upon the golden altar which is before the
throne of God. And the smoke of the incense of the prayers of the saints
ascended up before God from the hand of the angel.'' Andrew first reminds
his reader that the altar and the censer the saints see are shaped for them
in matter and colour, but are in themselves invisible and intelligible
realities. On such an altar the angel stands and offers the prayers of the
saints as incense to God. \gk{Thysiast\=erion de chrysoun ho Christos
estin; en h\=o pasa leitourgik\=e kai hagia synest\=eke dynamis}---the golden
altar is Christ, in whom every ministering and holy power stands together,
and on whom the sacrifices of the martyrs are offered. The throne before
which the altar stands is the order of the highest powers, so named because
of the fiery divine love poured into them and their pure wisdom and
knowledge, as the names of the powers nearest to God reveal: the Seraphim
are the burning, the Cherubim the fullness of knowledge. The angelic
offering of prayers is a priestly act done upon Christ the altar.

\paragraph{Michael and the dragon.} The thirty-fourth chapter treats the war
of the angels and the demons and the fall of Satan (Apoc 12:7--9).
Andrew applies the war to two moments: \gk{t\=e te pr\=ot\=e tou diabolou ex
alazoneias kai phthonou ekpt\=osei t\=es angelik\=es taxe\=os, t\=e te dia
tou Despotikou staurou kathairesei}, to the devil's first fall from the
angelic order through arrogance and envy, and to his overthrow by the
Lord's Cross, when ``the prince of this world is already judged'' (John 16:11) and
cast out of his ancient tyranny. The divine angels with their captain,
\gk{hama t\=o archistrat\=eg\=o Micha\=el}, with Michael their commander, not
bearing his arrogance, cast him from their company; so Ezechiel says that he
was driven by the Cherubim from the midst of the fiery stones, which Andrew
takes to be the angelic orders (Ezek 28:16). At Christ's coming the angels
who ministered to him after the temptation abhorred the devil again as a
disgraced slave. Andrew then records the Fathers' view of the time and
manner of the fall: after the making of the sensible world the devil was
cast down through pride and envy, having been entrusted first with the rule
of the air (Eph 2:2). He adds a sentence of Papias: to some of the angels
who were once divine God gave rule over the ordering of the earth and
charged them to rule well, and their order came to nothing (\emph{In
Apoc.}\ 34; PG~106). The Greek commentator thus agrees with Damascene that
the fallen prince was an angel set over the lower world, and with the whole
tradition that pride and envy were his sin.

\sectionguard
\subsection{Gregory Palamas: Image, Likeness, and Light}\label{sec:bz-palamas}

Gregory Palamas, archbishop of Thessalonica in the fourteenth century,
gathered his teaching into \emph{One Hundred and Fifty Chapters}, natural,
theological, moral, and practical. On the angels the chapters turn on two comparisons: the angels and man as images of God, and
the angels and man as partakers of the divine light. The Greek is
Migne's (PG~150), with a rendering after each passage.

\paragraph{Fellow-servants in the image.} The intelligible world is not to
be worshipped. \gk{Syndoulos oukoun h\=emin pasa noera physis, kai kat'
eikona tou ktisantos} (\emph{Cap.}\ 27; PG~150, 1140A): every intellectual
nature is our fellow-servant and made after the image of the Creator. The
angels are more honourable than we, \gk{h\=os s\=omat\=on ontes ektos, kai
t\=e as\=omat\=o pantapasi kai aktist\=o physei mallon engizontes}, being
outside bodies and nearer to the wholly bodiless and uncreated nature; those
who kept their order are our fellow-servants and yet honourable to us and far
above us in rank, while those who apostatized fell farthest from God and
became \gk{antitheoi}, opponents of God, and destroyers of our race.

\paragraph{Man more in the image.} In a series of chapters Palamas finds
the image of God more fully in man than in the angels. The rational soul
alone has mind, word, and life-giving spirit, and so \gk{mon\=e kai t\=on
as\=omat\=on angel\=on mallon kat' eikona tou Theou par' autou
ded\=emiourg\=etai; kai tout' echei ametapoi\=eton, kan m\=e epign\=o t\=en
heaut\=es axian} (\emph{Cap.}\ 39; PG~150, 1148B): it alone has been made by
God more after his image than even the bodiless angels, and this it keeps
unaltered even when it does not know its own dignity. Angel and soul are
both bodiless and not in a place, and neither is everywhere; but the soul
holds together and gives life to the body with which it was created, present
throughout it as containing it, and so bears in this too the image of God
(\emph{Cap.}\ 61). Then:

\begin{quote}
\gk{Ou kata touto monon mallon t\=on angel\=on ho anthr\=opos kat' eikona
pepoi\=etai Theou, hoti synektik\=en te kai z\=oopoion echei dynamin en
heaut\=o, alla kai kata to archein.} (\emph{Cap.}\ 62; PG~150, 1165A)
\end{quote}

\noindent Not in this alone is man made more after the image of God than the
angels, that he has in himself a power that holds together and gives life,
but also in ruling. The soul has a ruling part and a part that serves, and
God, because of what rules in us, gave man lordship over the whole earth.
The angels have no body yoked to them to be subjected to the mind; the
fallen hold their intellectual will fixed in evil, and the good hold theirs
fixed in good, \gk{h\=eniochou m\=edam\=os deomen\=en}, needing no
charioteer. The evil one did not possess earthly dominion but seized it; the
good angels were appointed by the Almighty to oversee the things of earth
after our fall and because of it; for, as Moses says in his song, God set
the bounds of the angels when he divided the nations (Deut 32:8), a
division Palamas dates after Cain and Seth. The next chapter adds
knowledge: the threefold character of our knowledge shows us more in God's
image than the angels, because we alone of creatures have sense joined to
intellect and reason, and so have discovered arts, sciences, and every kind
of knowledge---to till, to build, to bring forth from what is not, though
not from what is not at all, for that belongs to God alone; and God has
given to men alone that the invisible word of the mind should be heard in
the air, written, and seen through bodies, as a preparation for faith in the
coming of the Word through flesh, \gk{h\=on ouden oudam\=os metestin
angelois}, none of which belongs to angels at all (\emph{Cap.}\ 63; PG~150,
1165C--D).

\paragraph{The good angels more in the likeness.} The image is not the
likeness. \gk{All' ei kai to kat' eikona mallon h\=emeis t\=on angel\=on
echomen kai mechri nyn, pros to kath' homoi\=osin einai tou Theou poll\=o
elattoumetha, kai malista nyn, t\=on agath\=on angel\=on} (\emph{Cap.}\ 64;
PG~150, 1168A): even if we have the image more than the angels, even until
now, in being after the likeness of God we fall far short of the good angels,
and especially now. The perfection of the likeness is accomplished \gk{dia
t\=es ek Theou theias ellamps\=e\=os}, through the divine illumination from
God. The evil angels are deprived of it and are therefore under darkness;
the divine minds are filled with it and are therefore called a second light,
an outflowing of the first light. From this light the good angels have
knowledge even of sensible things, not by any sensing or natural power but
by a godlike power from which nothing present, past, or to come can be
hidden. The angels have less of the image and more of the likeness; man
has more of the image, and must receive the likeness through the same light.

\paragraph{The uncreated light of the angels.} Those who partake of this
illumination have it in measure, and know beings in proportion to their
measure. That the angels partake of it, that it is uncreated,
and that it is not the divine essence, Palamas says, all who read the
apostles and theologians know; and he calls on Dionysius: \gk{Kyklik\=os
gar, ph\=esin, hoi theioi kinoumenoi noes, henountai tais anarchois kai
ateleut\=etois ellampsesi tou kalou kai agathou} (\emph{Cap.}\ 65; PG~150,
1168C; \emph{DN} 4.8)---moving in a circle, he says, the divine minds are
united to the beginningless and endless illuminations of the Beautiful and
Good. Palamas draws two conclusions from the phrase. By speaking of
illuminations in the plural, Dionysius distinguished them from God's
essence, which is one; by calling them beginningless and endless, he showed
them to be uncreated. The angels live, in the Greek doctor's teaching, by
participation in the uncreated energies of God, which are God himself in his
self-giving and not his essence.

\paragraph{The Western counterpart.} Aquinas asks the same question of the
image and answers it with a distinction. ``as in their intellectual nature,
the angels are more to the image of God than man is, we must grant that,
absolutely speaking, the angels are more to the image of God than man is,
but that in some respects man is more like to God'' (\ST{I q.~93 a.~3}). The
respects are those Palamas names: in man ``a certain imitation of God''
is found, \latin{inquantum scilicet homo est de homine, sicut Deus de Deo;
et inquantum anima hominis est tota in toto corpore eius, et iterum tota in
qualibet parte ipsius, sicut Deus se habet ad mundum}---inasmuch as man is
from man as God is from God, and the soul is whole in the whole body and
whole in each part, as God is to the world. The Latin doctor calls the
angels more in the image \latin{simpliciter} and man \latin{secundum quid};
the Greek doctor calls man more in the image and the good angels more in the
likeness. On the light, the Latin schools teach that the blessed angels see
the divine essence itself by the created light of glory that raises their
intellect (\ST{I q.~12 a.~5}; \ST{I q.~62 a.~1}; see \ref{sec:q62}), and
Palamas teaches that angels and saints are united to God by the uncreated
illumination, which is distinct from the essence.

\sectionguard
\subsection{The Angels in the Divine Liturgy}\label{sec:bz-liturgy}

The Byzantine Liturgy, in the forms that bear the names of Chrysostom and
Basil, is built around the angels' presence. English quotations are from
Isabel F. Hapgood's \emph{Service Book of the Holy Orthodox-Catholic
Apostolic (Greco-Russian) Church} (1906), which follows the Russian use.

\paragraph{The entrance and the thrice-holy hymn.} At the Little Entrance,
when the Gospel book is carried into the sanctuary, the priest prays that
``with our entrance may enter also the holy Angels with us serving thee, and
with us glorifying thy goodness'' (p.~83; see \ref{sec:cm-ranks}). The choir
then sings the Trisagion, ``O Holy God, Holy Mighty, Holy Immortal One, have
mercy upon us'' (p.~86), while the priest says the prayer of the
thrice-holy hymn: ``O holy God, who restest in the Saints; who art hymned by
the Seraphim with a thrice-holy cry, and glorified by the Cherubim, and
adored by every heavenly Power \dots\ accept from the mouths of us sinners,
the Thrice-Holy song'' (p.~85). The Synaxarion of Constantinople relates how
the hymn was received. Under the emperor Theodosius and the patriarch
Proclus, when the city prayed outside the walls during the earthquakes, a
child was caught up from the midst of the people into the air; brought down
again, he cried with a great voice \gk{hoti hoi t\=on angel\=on choroi aneu
t\=es prosth\=ek\=es tou ``Ho staur\=otheis'' ton trisagion hymnon
anapempousi t\=o The\=o}, that the choirs of angels send up the thrice-holy
hymn to God without the addition ``who was crucified'': \gk{Hagios ho Theos,
hagios ischyros, hagios athanatos, ele\=eson h\=emas}; and at those words he
gave up his soul to God and the shaking of the earth ceased
(\emph{Synaxarium Constantinopolitanum}, 25 September, cols.\ 79--80). The
Church sings the Trisagion as the song she learned from the angels.

\paragraph{The Cherubic Hymn.} At the Great Entrance, when the gifts are
carried to the altar, the choir sings:

\begin{quote}
Let us, the Cherubim mystically representing, and unto the Life-giving
Trinity the thrice-holy chant intoning, all cares terrestrial now lay
aside. \dots\ That we may raise on high the King of all, like conqueror on
shield and spears, by the Angelic Hosts invisibly up-borne. (p.~94)
\end{quote}

\noindent The people represent the Cherubim; the King whom they receive is
borne up invisibly by the angelic hosts, as a victorious emperor was
carried on the shields of his soldiers. The priest's prayer during the hymn
confesses how great a thing it is to serve at this altar: ``For to serve
thee is a great and terrible thing even to the Heavenly Powers''; and he
addresses Christ as the one ``who art borne on the throne of the Cherubim;
who art Lord of the Seraphim and King over Israel'' (pp.~94--95).

\paragraph{The anaphora.} The eucharistic prayer gives thanks for the
ministry God accepts from men, ``although before thee stand thousands of
Archangels and myriads of Angels, with the Cherubim, and Seraphim,
six-winged, many-eyed, who soar aloft, borne on their pinions'' (p.~101;
Dan 7:10). In the anaphora of Basil every order is named: ``For Angels and
Archangels, Thrones, Dominions, Principalities, Authorities, Powers, and the
many-eyed Cherubim do laud thee. Before thee, round about, stand the
Seraphim, having each six wings'' (p.~102; Isa 6:2). The people sing the
Seraphim's hymn, ``Holy, holy, holy, Lord of Sabaoth,'' and the priest
continues, ``And we also, O Lord who lovest mankind, in company with these
blessed Powers do cry aloud and say: Holy art thou, and all-holy thou'' (p.~102).
Maximus's reading of the rite is the Liturgy's own: the faithful are
numbered among the angels by the thrice-holy hymn before they are brought to
the Father. The Latin Mass says the same in its prefaces, \latin{cum Angelis
et Archangelis, cum Thronis et Dominationibus}, and in the Canon's prayer
that an angel carry the offering to the altar on high (see
\ref{sec:cm-ranks}).

\sectionguard
\subsection{The Synaxis of the Bodiless Powers and the Miracle at Chonae}\label{sec:bz-feasts}

\paragraph{The Synaxis of 8 November.} The Byzantine calendar keeps on 8
November ``The Feast of St.\ Michael the Archangel, and of the other
Bodiless Powers of Heaven'' (Hapgood, \emph{Service Book}, p.~xiv, where it
is dated Nov.\ 8--21, old and new style). The Synaxarion of Constantinople
gives the reason of the name. When the all-good and man-loving God willed to
make this world, he first conceived the angelic and heavenly powers, and
the conception became work; he set over each order a taxiarch and captain.
One of the taxiarchs grew proud, thought himself a god, fell from his rank,
and was cast out of heaven with the order under him. Then Michael, captain
of another order, seeing the apostate fallen, gathered the choirs of the
angels and, saying \gk{Prosch\=omen}, let us attend, hymned the God of
all, as if to say: let us attend, we who were made creatures and stand
before God; what has befallen those who until now were light with us and
now have become darkness? \gk{H\=e toiaut\=e oun synkrot\=esis \=onomasth\=e
synaxis t\=on as\=omat\=on, tout' estin prosoch\=e kai homonoia kai
hen\=osis}: such a gathering was named the synaxis of the bodiless, that is,
attention and concord and union (\emph{Synaxarium Constantinopolitanum}, 8
November, cols.\ 203--204). The feast is the angels' own liturgy of
fidelity, kept on the day their unity was sealed by Michael's word. The
Synaxarion names several churches of the capital where it was celebrated.
Theodore's sermon, above, was preached for it.

The Russian service book appoints for the service of the Bodiless Powers
the epistle Heb 2:2--10 and the gospel Luke 10:16--21 (p.~xxiv), and it
gives among the hymns of the weekdays the hymn of Monday, which the Church
dedicates to the angels: ``O ye Chieftains of God, servitors of the divine
glory, Guides of men and Captains of the Bodiless Powers: Entreat that which
is profitable for us, and great mercy; in that ye are the Chieftains of the
Bodiless Powers'' (p.~61).

\paragraph{The miracle at Chonae.} On 6 September the same Synaxarion
keeps the feast of the archangel Michael, \gk{anamn\=esin poioumenoi tou
thaumatos}, making memory of the wonder he worked at Colossae in Phrygia,
afterward called Chonae (cols.\ 19--20). Tradition relates that John the
Theologian, preaching at a place near Colossae called Chaeretopa, foretold
an appearing of the archangel there and the rising of a healing spring;
when it came to pass, many were healed by its water, and the believers built
a beautiful church. A certain Archippus served in it for seventy years.
Moved by envy, the pagans first tried to turn a river into the holy water,
but the river was turned aside; then they dug a great trench, gathered two
rivers behind a dam, and loosed them at midnight to sweep away Archippus,
the church, and the spring. Archippus prayed with tears to the archangel,
and Michael came, appearing as a pillar of fire, bade him not fear, stood on
a great rock toward which the waters were rushing, struck and split it with
his staff, and commanded the waters of the two rivers \gk{ch\=oneuesthai},
to be funnelled into it. From that wonder the place has been called Chonae.
For these reasons, the Synaxarion concludes, the Church of Constantinople
gathers at the archangel's shrine at Anaplus to keep his
memory. The Western Church keeps the dedication of Michael's church on 29
September, and the Missal of 1962 gives it the first class, praying \latin{Deus,
qui, miro ordine, Angelorum ministeria hominumque dispensas}, O God, who
dispose in wonderful order the ministries of angels and of men, that those
who always minister before God in heaven may protect our life on earth
(\emph{Missale Romanum}, 1962, 29 September, n.~3727).

\sectionguard
\subsection{Witnesses of the Byzantine Line}\label{sec:bz-table}

\begin{fourstable}{Witness}{Locus}{Teaching}{Western counterpart}
Damascene & \emph{De fide orth.}\ II.3 (PG~94, 865) & The angel is an intelligent essence, ever-moving, free, incorporeal, ministering to God, immortal by grace & \ST{I q.~50 a.~1 ad~1--2}; a.~5 ad~1\\
Damascene & II.3 & Not susceptible of repentance because incorporeal; the good now immovable by grace & \ST{I q.~62 a.~8}; \ST{I q.~64 a.~2}\\
Damascene & II.3 & In heaven they are not on earth; present in mental places, not in several at once & \ST{I q.~52 a.~2} (s.c.)\\
Damascene & II.3 & Equality of essence unknown; difference in brightness and position; higher illumine lower & \ST{I q.~50 a.~4} (each angel a species); \ST{I q.~106 a.~1}\\
Damascene & II.3 (PG~94, 873--874) & Nine orders in three triads, after Dionysius & \ST{I q.~108 a.~6} (Dionysius and Gregory compared)\\
Damascene & II.3 & With Gregory the Theologian, angels created before the visible world & \ST{I q.~61 a.~3} (with the corporeal world, the more probable; the contrary not erroneous)\\
Damascene & II.4 (PG~94, 877--878) & The devil, set over the earth, fell by free choice; the fall is to angels what death is to men & \ST{I q.~63 a.~7}; \ST{I q.~64 a.~2}; \ST{I q.~110 a.~1}\\
Damascene & \emph{De imag.}\ III & Angels may be depicted and circumscribed after their own nature; they are honoured as God's ministers & Nicaea II (DS~600--601)\\
Maximus & \emph{Myst.}\ 13, 19, 24 (PG~91, 692, 696, 709) & The thrice-holy hymn numbers the faithful with the angels and foreshadows equal honour with them & \emph{Missale Romanum}, prefaces (\latin{cum Angelis et Archangelis}); \ST{I q.~108 a.~8}\\
Nicaea II & DS~600--601 & Images of the angels are set up and honoured, not with latria & Constantinople IV, can.~3 (DS~656)\\
Theodore the Studite (as edited by Mai) & \emph{Or.}\ 6.1--3 (PG~99, 729--736) & Michael leads the choir above and below; one essence, many offices; almost every mystery of God on earth is done through angels & \ST{I q.~110 a.~1}; \ST{I q.~113 aa.~1--4}\\
Theodore the Studite & \emph{Or.}\ 5.5 (PG~99, 728) & All nine orders attend the Dormition & \emph{Missale Romanum}, 15 August (\latin{gaudet exercitus Angelorum})\\
Andrew of Caesarea & \emph{In Apoc.}\ 21 (PG~106) & The angelic powers offer the prayers of the saints on Christ the altar & Roman Canon, \latin{Supplices te rogamus}; \ST{III q.~83 a.~4 ad~9}\\
Andrew of Caesarea & \emph{In Apoc.}\ 34 (PG~106) & The devil fell through arrogance and envy, having ruled the air; Michael cast him out & \ST{I q.~63 aa.~2, 7}\\
Palamas & \emph{Cap.}\ 27, 39, 61--63 (PG~150, 1140--1165) & Man more in the image than the angels, as life-giver of a body, ruler, and knower through sense & \ST{I q.~93 a.~3} (angels more in the image absolutely; man in some respects)\\
Palamas & \emph{Cap.}\ 64--65 (PG~150, 1168) & The good angels more in the likeness; they partake of the uncreated illumination, distinct from the essence & \ST{I q.~12 a.~5}; \ST{I q.~62 a.~1} (the created light by which the intellect sees God's essence)\\
Divine Liturgy & Hapgood 1906, pp.~83--102 & The angels enter with the clergy; the people represent the Cherubim; the anaphora joins the Seraphim's hymn & Roman prefaces and Sanctus (\emph{Missale Romanum}, 1962)\\
Synaxarion of Constantinople & 25 Sept.\ (cols.\ 79--80) & The Trisagion learned from the angels' choirs & ---\\
Synaxarion of Constantinople & 8 Nov.\ (cols.\ 203--204) & The synaxis of the bodiless: Michael's \gk{Prosch\=omen} after the fall & 29 September (\emph{Missale Romanum}, 1962)\\
Synaxarion of Constantinople & 6 Sept.\ (cols.\ 19--20) & Michael's miracle at Chonae & 29 September\\
\end{fourstable}
